\scanpage{38}
% L1–4, c)
c) Supposons les $\varphi^n$ des isom \deleted{\unsure{pour $i\le n$|fonctoriels}}, \deleted{\unsure{un épimorphisme|un épimorphisme surjectif}} Quand peut-on conclure que $\varphi^{n+1}$ est un mono ? Réponse : il suffit que \deleted{\unsure{$\varphi$}} $T^{n+1}$ soit effaçable. Donc si les $T^i$ \unsure{son|sont} effaçables pour $i>i_0$, et les $\varphi^i$ des \unsure{d° \upshape isom|isom} pour $i$ \deleted{\unsure{petit}} $<i_0$, un

% L5–7, crossed block
\deleted{d) Supposons les $\varphi^n$ des isom, pour \unsure{$i\le n$}, un mono, en \unsure{dim $n+1$}, Quand peut-on affirmer que $\varphi^{n+1}$ est \unsure{$\overline{\mathrm M}$} un isom ? \deleted{Quand} \inserted{Il suffit que} $T'{}^{n+1}$ \deleted{\unsure{$(B)$}} soit coeffaçable}

% L8–10, resumption of c)
mono pour $i\unsure{>|\ge}i_0$, et $\varphi^i$ un épim pour $i\le n$ (où $n\ge i_0$) alors $\varphi^i$ est un \deleted{\unsure{épim|épi}} \uline{isom} pour $i\le n$, un \unsure{épim|épi} en dim $n+1$.

% L11–13
Revenons sur a), b). On a des propriétés \deleted{\unsure{$I$}} $I_n$, $J_n$ d'un argument $A\in\mathrm{Ob}\,C$,
\[
\left(\begin{aligned}
 I_n(A)&:\varphi^{\deleted{n}\inserted{i}}(A)\ \text{isom si }i\le n\\
 J_n(A)&:\varphi^i(A)\ \text{isom si }i\le n,\ \text{mono}\\
 &\hspace{7em}\text{si }i=n+1
\end{aligned}\right]
\]

% L14–16
\[
\left\{\begin{aligned}
 J_n&\Rightarrow I_n\Rightarrow J_{n-1}\quad\ldots\\
 I_n\ \text{pour tt }A,\ +J_n\ \text{pour cog.}&\to J_n\ \text{pour tout } A\\
 J_n\ \text{pour tt }A,\ +I_{n+1}\ \text{pour tt cog}&\Rightarrow I_{n+1}\ \text{pour tt }A
\end{aligned}\right.
\]

% L17–20
\uline{\unsure{Concl|Concl.}} Pour vérifier $I_n$ (resp $J_n$) pour tt $A$, sachant que c'est vrai pour $n$ petits, il suffit de vérifier $I_m(B)$ (resp $J_m(B)$) pour $m\le n$ et tt cog. $B$.

% L21, double rule
\rule{3cm}{0.4pt}\\[-1.5ex]\rule{3cm}{0.4pt}

% L22–27
On avait ici des arguments où on ``\unsure{montait|remontait}'' les degrés. Dans les questions de dim \unsure{ch.|coh.}, au contraire, il faut ``descendre''. La question est de donner des critères pour qu'un $T^i$ soit nul (pour tt \uline{argument}). Ici on doit utiliser les faits suivants (\unsure{tenant déjà|tenant d'abord|savoir aussi dans} au fait que $T^i$ est

% archive sheet number: 37
